% ============================================================================
% NSF ENG CMMI, solicitation NSF 26-515 (posted Aug 17, 2026)
% One-page project summary for the Advanced Manufacturing, Design, and
% Operations (AMDO) program, to email to cmmi-amdo@nsf.gov before submitting.
% Solicitation, Sec. II: "Proposers are strongly encouraged to email a
% one-page project summary to the Cognizant Program Director(s) before
% submitting a full proposal for guidance on whether the proposed research
% topic falls within the scope of the program."
% The solicitation sets no format for this summary; we use Times New Roman
% 11 pt and 1 in margins. Sources and open questions are in README.md.
%
% Base text: afosr-yip-white-paper/afosr-yip-white-paper-a2c.tex.
% Simulation numbers are only those in that paper's Section 5.
%
% Build:  xelatex one-page-summary.tex
% Lint:   cringe-filter lint  -c proposal one-page-summary.tex
%         cringe-filter score -c proposal one-page-summary.tex
% ============================================================================
\documentclass[11pt,letterpaper]{article}

\usepackage[margin=1in]{geometry}
\usepackage{fontspec}
\IfFontExistsTF{Times New Roman}
  {\setmainfont{Times New Roman}}
  {\setmainfont{TeX Gyre Termes}}
\usepackage{microtype}
\usepackage{xcolor}
\usepackage[colorlinks=true,urlcolor=blue!70!black]{hyperref}

\setlength{\parindent}{0pt}
\setlength{\parskip}{3pt plus 1pt minus 1pt}
\pagestyle{empty}

\newcommand{\head}[1]{\textbf{#1.}\ }
\hyphenation{re-POW-DER AMA-ZE-MET}

\begin{document}

\begin{center}
{\large\bfseries Choosing the Next Experiment in a Multi-Step Powder-to-Part
Workflow,\\ with Printable High-Strength Aluminum Alloys as the Test
Case\par}
\vspace{4pt}
Sterling G.\ Baird, Ph.D., Assistant Professor of Mechanical Engineering\\
Brigham Young University (BYU), Provo, UT 84602;
\href{mailto:sterling.baird@byu.edu}{sterling.baird@byu.edu}; (801) 422-5344\\
Summary for the Advanced Manufacturing, Design, and Operations (AMDO)
program, NSF 26-515
\end{center}
\vspace{-4pt}

\head{Research question}
Developing a new alloy for laser powder bed fusion (LPBF) runs through a chain
of steps, from a thermodynamic calculation to powder, printed tracks and
cubes, microscopy, and tensile tests. A calculation takes minutes and a
tensile test takes days. Campaigns often send every candidate through the
whole chain, or screen many by simulation and test a few. We ask how to
decide, as data arrive, which step of a multi-step manufacturing workflow to
run next and on which candidate, while carrying uncertainty through the
chain. Two features make this harder: several properties compete, and results
from many experiments in progress return at different times. Our test case is high-strength
aluminum, since alloys such as 7075 crack as they solidify during LPBF.
Grain-refining additions such as zirconium can prevent the cracking (Martin et
al., \emph{Nature} 2017), but each change in composition means the laser
settings must be tuned again.

\head{Approach}
We will give each step its own statistical model and link the models in
order, so the outputs of one step feed the next. Recent work on
such linked models picks the experiment with the largest expected gain per
unit cost and lets a candidate stop at any step (Astudillo and Frazier,
NeurIPS 2021; Buathong et al., ICML 2024). It has been shown for one
objective and one experiment at a time, on test functions and simulators. We
will extend it to weigh competing objectives (freedom from cracks and pores,
strength, and ductility) with no weighting known in advance, and to choose
new experiments while others are pending, since microscopy results arrive days
after printing. We will demonstrate it on the five-step workflow,
with real build-to-build scatter and with incomplete physics inputs from
calculation of phase diagrams (CALPHAD), which ignores the fast cooling of the
melt pool. We will test whether the method reaches a target set of trade-offs at
lower total cost than standard multi-objective Bayesian optimization run on
the full chain. We will also test whether its predicted chance of a dense,
crack-free build matches observed rates on held-out builds.

\head{Resources and preliminary results}
Our laboratory atomizes aluminum alloy powder (AMAZEMET rePOWDER
ultrasonic atomizer, induction melting), four compositions per week. We buy
LPBF builds from a University of Utah user facility (\$800 per build) and use
the BYU Microscopy Center. Since the
lab started in December 2025, we have run CALPHAD on BYU high-performance
computing and completed a powder-to-part demonstration on a baseline aluminum
alloy, with microscopy and tensile data. In a hand-built, noise-free
simulation with one objective, where the CALPHAD crack screen missed
zirconium's effect, screening and then tuning stalled at 63\% of the optimum
after \$100,000. The cost-aware search reached that level after about
\$26,500 (3.8 times less spend; 90\% confidence interval 2.0 to 5.7).
Screening kept pace when its physics was right.

\head{Intellectual merit}
To our knowledge, neither a multi-objective nor an asynchronous version of
cost-aware optimization over linked steps has been developed, and none has
been tested on a physical manufacturing workflow. The project will show
how to carry uncertainty through a multi-step process while keeping model
uncertainty, which more data reduces, separate from build-to-build scatter,
which it does not. The decision method is the subject of the research and
applies to other multi-step workflows, such as qualifying a new material or
process.

\head{Broader impacts}
Undergraduates will make powder and prepare specimens, and a Ph.D.\ student
will run the printing and microscopy campaigns, all gaining training in metal
additive manufacturing and autonomous experimentation. We will release
software, data, and models openly each year so other laboratories can apply
the method to their own materials, printers, and tests.

\head{Question for the program director}
Does this topic fall within the scope of AMDO, or would it fit better in
another program under NSF 26-515?

\end{document}
